\section{Super-heavies}

\subsection{AX1-0 Tiger Shark}

The Tiger Shark is the Barracuda's larger cousin. It is an enormous aircraft deployed in the role of fighter-bomber. Tiger Sharks are never as numerous as Barracudas, but they have many characteristics and systems in common with them. All crew members come from the Air Caste, which gives them the natural advantage of a superior three-dimensional awareness and a tolerance for higher acceleration speeds, and more G in the turn, than a human pilot. Just like the Barracuda, the different Air Castes of the Septs operate according to slightly different versions.

The Tiger Shark has the \textbf{Wounds (2)}, \textbf{Integral Armour}, and \textbf{Flyer} abilities. Its Ion Cannons have \textbf{Damage (+1)}. Its Burst Cannons have the \textbf{Turret} ability.

\begin{unitstatstable}
\SetCell[r=2]{c,m} -- &
\SetCell[r=2]{c,m} 3+ &
\SetCell[r=2]{c,m} +3 &
\SetCell[r=2]{c,m} 5 &
Heavy Ion Cannons &
45 cm &
2 &
4+ &
$-3$
\\
& & & &
Burst Cannons &
20 cm &
4 &
4+ &
0
\end{unitstatstable}

\unitimage{tiger-shark-ax1-0.png}

\unitdivider

\subsection{Scorpionfish}

After facing Imperial super-heavy tanks during the Damocles Crusade and the conflicts that followed, the T'au devoted considerable effort to developing their own super-heavy tanks. Having no other vehicle large enough to serve as a base, the Orca transport was heavily armed and armoured, reducing its mobility so much that it became a surface antigrav. This variant was called the Scorpionfish Heavy Missile Tank and carries no less than a complete set of submunition, seeker, and tracer missiles as well as a direct-fire missile launcher. It does not have a single main weapon but rather acts as a platform for many smaller systems, making it capable of attacking any kind of target. This flexibility has proved a considerable asset in recent T'au campaigns. It forms part of the rear firing line of T'au armies, protecting the rear lines as solidly as a rock.

The Scorpionfish has the \textbf{Heavy Antigrav}, \textbf{Wounds (2)}, and \textbf{Markerlight (75 cm)} abilities. It possesses a Smart Missile system. Its Heavy Missile Launcher has three firing modes which all use \textbf{Seeker Missiles}. The Submunition mode uses a \textbf{Template (12 cm)} that \textbf{Reduces Cover ($-1$)}.

\begin{unitstatstable}
\SetCell[r=5]{c,m} 15 cm &
\SetCell[r=5]{c,m} 2+ &
\SetCell[r=5]{c,m} +2 &
\SetCell[r=5]{c,m} 5 &
Smart Missiles &
45 cm &
2 &
4+ &
$-1$
\\
& & & &
\SetCell[c=5]{c,m}\textit{Heavy Missile Launcher (choose)} &
& & &
\\
& & & &
Tracer Missiles &
100 cm &
3 &
3+ &
$-3$
\\
& & & &
Submunition Missiles &
100 cm &
Template &
4+ &
0
\\
& & & &
Seeker Missiles &
100 cm &
6 &
4+ &
$-1$
\end{unitstatstable}

\unitimage{scorpionfish.png}

\unitdivider

\subsection{Tiger Shark Drone Transport}

The Tiger Shark plays both the role of a Drone transport and a heavy support fighter. Its lower hold can transport 4 Attack Drone bases, which can be dropped at altitude to reinforce T'au lines or to launch surprise attacks behind enemy lines. The Tiger Shark does not have to land to disembark its Drones and may drop Attack Drone bases at any point during its movement.

The Tiger Shark has the \textbf{Wounds (2)}, \textbf{Integral Armour}, \textbf{Flyer}, \textbf{Tiger Shark}, and \textbf{Transport (4 Drones)} abilities. Its Burst Cannons have the \textbf{Turret} ability.

\begin{unitstatstable}
\SetCell[r=2]{c,m} -- &
\SetCell[r=2]{c,m} 3+ &
\SetCell[r=2]{c,m} +1 &
\SetCell[r=2]{c,m} 5 &
Ion Cannons &
30 cm &
2 &
4+ &
$-2$
\\
& & & &
Burst Cannons &
20 cm &
4 &
4+ &
0
\end{unitstatstable}

\unitimage{tiger-shark-drone-transport.png}

\unitdivider

\subsection{Orca Transport}

The Orca is an orbital transport vehicle specially designed to carry troops, equipment, and supplies from spacecraft in orbit to a planet's surface. Unlike the Thunderhawk gunship of the Space Marines or the much larger T'au Manta, the Orca is not a front-line combat vehicle, armed and armoured to take part in a battle.

The Orca is equipped only for self-defence. Its primary function is transport. The Orca has the \textbf{Integral Armour}, \textbf{Wounds (2)}, \textbf{Transport (8)}, and \textbf{Flyer} abilities. Its Burst Cannons have the \textbf{Turret} ability.

\begin{unitstatstable}
\SetCell[r=2]{c,m} 30 cm &
\SetCell[r=2]{c,m} 3+ &
\SetCell[r=2]{c,m} +1 &
\SetCell[r=2]{c,m} 5 &
Burst Cannons &
20 cm &
4 &
4+ &
0
\\
& & & &
Smart Missiles &
30 cm &
2 &
4+ &
$-1$
\end{unitstatstable}

\unitimage{orca.png}

\unitdivider

\subsection{Swordfish}

Another result of T'au work in the design of super-heavy vehicles, the Swordfish super-heavy tank is the T'au answer to Imperial Shadowswords and Baneblades. Even if they had not initially intended to use super-heavy vehicles (except the Manta), the T'au quickly realised the effectiveness of Imperial super-heavies and rapidly developed their own. The first to be built was the Scorpionfish and shortly after came the Swordfish.

The Swordfish has the \textbf{Heavy Antigrav} and \textbf{Wounds (2)} abilities. Its Heavy Railgun has the \textbf{Turret} and \textbf{Damage (+1)} abilities. Finally, it possesses a Smart Missile system.

\begin{unitstatstable}
\SetCell[r=3]{c,m} 15 cm &
\SetCell[r=3]{c,m} 2+ &
\SetCell[r=3]{c,m} +2 &
\SetCell[r=3]{c,m} 5 &
Twin Heavy Railgun &
90 cm &
2 &
4+ &
$-3$
\\
& & & &
Smart Missiles &
45 cm &
2 &
4+ &
$-1$
\\
& & & &
Burst Cannons &
30 cm &
2 &
4+ &
0
\end{unitstatstable}

\unitimage{swordfish.png}
